\section{计算机系统概述}

\subsection{计算机发展历程}

\begin{mydef}{冯·诺依曼结构}{von-neumann}
冯·诺依曼（John von Neumann）于 1945 年提出「存储程序」（Stored Program）概念，奠定了现代通用计算机的基本架构。其核心思想可概括为：

\begin{enumerate}
    \item \textbf{存储程序}：指令和数据以同等地位存放在同一个存储器中，按地址访问。
    \item \textbf{顺序执行}：CPU 从存储器中逐条取出指令，经译码后执行，程序计数器（PC）自动指向下一条指令。
    \item \textbf{二进制编码}：指令和数据均以二进制形式表示。
    \item \textbf{五大部件}：计算机由\textbf{运算器}、\textbf{控制器}、\textbf{存储器}、\textbf{输入设备}和\textbf{输出设备}五大部分组成。
\end{enumerate}
\end{mydef}

\begin{conclusion}{冯·诺依曼结构 vs 现代计算机结构}{vn-vs-modern}
\textbf{冯·诺依曼结构的特点：以运算器为中心。}

在冯·诺依曼原始设计中，运算器（ALU）处于整个系统的中心位置：输入/输出设备与存储器之间的数据传送都需要经过运算器，运算器既是计算的核心，也是数据流动的枢纽。这一设计导致了所谓的「冯·诺依曼瓶颈」——总线上的数据传送成为性能制约因素。

\textbf{现代计算机结构的特点：以存储器为中心。}

现代计算机体系结构转向了「以存储器为中心」：
\begin{itemize}
    \item CPU（运算器 + 控制器）与存储器通过系统总线直接通信；
    \item I/O 设备通过 I/O 控制器/接口电路与总线相连，数据可以在存储器和 I/O 设备之间直接传送（DMA 方式），无需经过 CPU 的运算器；
    \item 存储器成为数据交换的中心枢纽。
\end{itemize}
\end{conclusion}

\begin{attention}
\textbf{408 常考点：} 冯·诺依曼结构的核心特征是「存储程序」。多选题中常出现「冯·诺依曼计算机的特点」的判断，注意：
\begin{itemize}
    \item 「指令和数据以同等地位存放在同一存储器」——正确（但现代机器中指令 Cache 和数据 Cache 分离是性能优化，不改变「同一存储器」的逻辑模型）；
    \item 「指令与数据均用二进制表示」——正确；
    \item 「指令顺序执行，PC 自增指向下一条」——正确（是基本模型；现代 CPU 的流水线和乱序执行是内部优化，程序员看到的是顺序执行）。
\end{itemize}
\end{attention}

\begin{conclusion}{计算机发展简史}{computer-history}
\begin{tablebox}{计算机发展的四个时代}
\centering
\begin{tabular}{|c|c|c|l|}
\hline
时代 & 时间 & 核心元件 & 特点 \\
\hline
第一代 & 1946--1957 & 电子管 & 体积大、功耗高、机器语言编程 \\
第二代 & 1958--1964 & 晶体管 & 出现高级语言（FORTRAN）、批处理系统 \\
第三代 & 1965--1971 & 中小规模 IC & 操作系统成熟、出现微机雏形 \\
第四代 & 1972--至今 & 大规模/超大规模 IC & 微处理器诞生、PC 普及、多核并行 \\
\hline
\end{tabular}
\end{tablebox}

\textbf{摩尔定律（Moore's Law）：} 集成电路上可容纳的晶体管数目约每 18--24 个月翻一番，性能也随之提升。这一定律至今已接近物理极限，现代计算转向多核并行和专用加速器（GPU/TPU）方向。
\end{conclusion}

\subsection{计算机硬件组成}

\begin{mydef}{计算机硬件五大部件}{hardware-components}
现代计算机的硬件系统由以下五大部件组成，通过系统总线（地址总线、数据总线、控制总线）互联：

\begin{enumerate}
    \item \textbf{运算器（ALU，Arithmetic Logic Unit）}：执行算术运算和逻辑运算。核心部件包括 ALU、累加器（ACC）、乘商寄存器（MQ）、通用寄存器组等。
    \item \textbf{控制器（CU，Control Unit）}：指挥整个计算机系统协调工作。核心部件包括程序计数器（PC）、指令寄存器（IR）、指令译码器（ID）、时序发生器和微操作信号发生器。运算器与控制器合称 CPU。
    \item \textbf{存储器（Memory）}：存放程序和数据。分为主存储器（内存）和辅助存储器（外存）。主存由存储体、MAR（存储器地址寄存器）、MDR（存储器数据寄存器）及读写电路组成。
    \item \textbf{输入设备}：将程序和数据从外界送入计算机，如键盘、鼠标、扫描仪。
    \item \textbf{输出设备}：将处理结果从计算机送出到外界，如显示器、打印机。
\end{enumerate}
\end{mydef}

\begin{attention}
\textbf{408 常考辨析：CPU 包含哪两部分？} 运算器 + 控制器 = CPU（中央处理器）。注意「主机」的概念在 408 中通常指 CPU + 主存储器（不包括外存和 I/O 设备）。
\end{attention}

\begin{conclusion}{运算器的核心寄存器}{alu-registers}
\begin{tablebox}{运算器常用寄存器}
\centering
\begin{tabular}{|c|l|}
\hline
寄存器 & 功能 \\
\hline
ACC（Accumulator） & 累加器，存放 ALU 运算结果 \\
MQ（Multiplier-Quotient） & 乘商寄存器，乘法时存乘数，除法时存商 \\
X（通用操作数寄存器） & 存放一个 ALU 输入操作数 \\
PSW/FLAGS & 程序状态字/标志寄存器，存放条件码（ZF/OF/CF/SF）\\
\hline
\end{tabular}
\end{tablebox}

不同体系结构中寄存器命名有所不同，但功能等价。x86 中累加器为 EAX/RAX，ARM 中使用通用寄存器 r0--r15。
\end{conclusion}

\begin{conclusion}{控制器的基本工作流程}{cu-workflow}
控制器完成一条指令的基本流程（取指--译码--执行）：
\begin{enumerate}
    \item \textbf{取指（Fetch）}：按 PC 所指地址从存储器取指令$\to$IR，PC 自增指向下一条指令。
    \item \textbf{译码（Decode）}：指令译码器分析 IR 中的操作码，产生相应的控制信号。
    \item \textbf{执行（Execute）}：根据译码结果，控制相关部件完成数据传送或运算。
    \item \textbf{写回（Write Back）}：将结果写回寄存器或存储器（如有需要）。
\end{enumerate}
\end{conclusion}

\begin{conclusion}{存储器的层次结构}{memory-hierarchy}
存储器按速度和容量分为层次结构（见后续章节详述）：
\[
\text{CPU 寄存器} > \text{Cache（L1/L2/L3）} > \text{主存（DRAM）} > \text{辅存（SSD/HDD）}
\]
从左到右：速度递减、容量递增、单位成本递减。
\end{conclusion}

\subsection{计算机软件系统}

\begin{mydef}{系统软件 vs 应用软件}{software-classification}
计算机软件分为两大类：

\begin{enumerate}
    \item \textbf{系统软件（System Software）}：管理计算机系统资源、支撑其他软件运行的软件。包括：
    \begin{itemize}
        \item \textbf{操作系统（OS）}：管理硬件资源、提供程序运行环境（如 Linux、Windows、macOS）。
        \item \textbf{语言处理程序}：编译器、汇编器、解释器，将高级/汇编语言程序翻译为机器码。
        \item \textbf{数据库管理系统（DBMS）}：如 MySQL、PostgreSQL。
        \item \textbf{实用程序}：加载器、链接器、调试器。
    \end{itemize}
    \item \textbf{应用软件（Application Software）}：直接为用户完成特定任务的软件，如办公软件、浏览器、游戏、科学计算程序等。
\end{enumerate}

\textbf{408 常考辨析：} 操作系统属于系统软件；数据库管理系统（DBMS）通常也归入系统软件范畴；编译器、汇编器是系统软件。
\end{mydef}

\begin{conclusion}{三个级别的程序设计语言}{language-levels}
\begin{tablebox}{三种语言级别对比}
\centering
\begin{tabular}{|c|c|c|c|}
\hline
特性 & 机器语言 & 汇编语言 & 高级语言 \\
\hline
形式 & 二进制代码 & 助记符 & 类自然语言/数学表达式 \\
可读性 & 极差 & 较差 & 良好 \\
可移植性 & 无 & 差 & 好（需编译器适配）\\
执行效率 & 最高 & 高 & 取决于编译器优化 \\
程序示例 & \texttt{B8 0001} & \texttt{MOV AX, 1} & \texttt{int a = 1;} \\
翻译方式 & 直接执行 & 汇编器$\to$机器码 & 编译器$\to$机器码或解释执行 \\
\hline
\end{tabular}
\end{tablebox}

\textbf{翻译过程：} 高级语言 $\xrightarrow{\text{编译器/解释器}}$ 汇编语言 $\xrightarrow{\text{汇编器}}$ 机器语言。注意：编译与解释是两种不同的翻译方式：
\begin{itemize}
    \item \textbf{编译（Compilation）}：一次性将整个源程序翻译为目标代码，再执行目标代码（如 C/C++、Go、Rust）。
    \item \textbf{解释（Interpretation）}：逐条翻译并执行，不生成独立的目标文件（如 Python、JavaScript 的传统执行方式）。
\end{itemize}
\end{conclusion}

\subsection{计算机系统的层次结构}

\begin{mydef}{计算机系统的多层次模型}{system-layers}
计算机系统可以抽象为层次模型：下层为上层提供接口（服务），上层通过调用下层接口实现更高级的功能。从底层到上层的典型分层如下：
\end{mydef}

\begin{conclusion}{从微程序到高级语言的六级层次}{system-layers-detail}
\begin{center}
\begin{tabular}{|c|c|l|}
\hline
层次 & 层级 & 核心概念 \\
\hline
$L_0$ & 微程序级/硬联逻辑 & 微指令、微操作，由硬件直接执行 \\
$L_1$ & 传统机器语言级 & 二进制机器指令，由微程序解释执行 \\
$L_2$ & 操作系统级 & 提供系统调用，管理资源与抽象 \\
$L_3$ & 汇编语言级 & 助记符指令，由汇编器翻译为 $L_1$ \\
$L_4$ & 高级语言级 & C/C++/Java 等，由编译器翻译 \\
$L_5$ & 应用层（用户层）& 应用程序通过 API/库使用系统服务 \\
\hline
\end{tabular}
\end{center}

\textbf{关键概念：}
\begin{itemize}
    \item \textbf{翻译（Translation）}：将上层程序整体转换为下层等价程序，再在下层执行。$L_3\to L_1$（汇编器）、$L_4\to L_1$（编译器）均通过翻译。
    \item \textbf{解释（Interpretation）}：下层程序逐条处理上层程序的每条语句并直接执行，不生成完整的目标程序。$L_1$ 的机器指令通常由 $L_0$ 的微程序解释执行。
    \item \textbf{虚拟机（Virtual Machine）}：通过「翻译 + 解释」的组合，使得上层语言不需要考虑底层的硬件细节，每一层对上层来说都是一台「虚拟机器」。典型例子：JVM 将 Java 字节码解释/即时编译（JIT）在物理机上执行。
\end{itemize}
\end{conclusion}

\begin{attention}
\textbf{408 常考点——翻译 vs 解释：}
\begin{itemize}
    \item 在层次结构模型中，$L_3$（汇编）到 $L_1$（机器）是\textbf{翻译}；$L_4$（高级语言）到 $L_1$ 也是\textbf{翻译}。
    \item $L_0$（微程序）到 $L_1$ 是\textbf{解释}——每一条机器指令由一段微程序解释执行。
    \item 注意：物理上不存在 $L_2$--$L_5$ 层的「硬件」，这些是虚拟机器（由下层软件+硬件模拟出来的抽象）。物理机器一般指 $L_0$ 和 $L_1$。
\end{itemize}
\end{attention}

\subsection{计算机性能指标}

\begin{mydef}{基本性能指标}{performance-metrics}
衡量计算机性能的主要指标包括：
\begin{enumerate}
    \item \textbf{机器字长（Word Size）}：CPU 一次能处理的二进制数据的位数。通常与通用寄存器位数和 ALU 宽度一致。字长越长，计算精度越高，但硬件成本也越大。常见：32 位、64 位。
    \item \textbf{主频（Clock Frequency）}：CPU 内部时钟信号的频率，以 Hz 为单位。主频 $f$ 决定了时钟周期 $T = 1/f$。例如主频 3.0 GHz $\to$ 时钟周期 $T \approx 0.33\text{ ns}$。
    \item \textbf{CPI（Cycles Per Instruction）}：执行一条指令所需的平均时钟周期数。CPI 越小，CPU 效率越高。
    \item \textbf{MIPS（Million Instructions Per Second）}：每秒执行百万条指令。
    \item \textbf{MFLOPS（Million Floating-point Operations Per Second）}：每秒百万次浮点运算。类似还有 GFLOPS、TFLOPS。
\end{enumerate}
\end{mydef}

\begin{formula}{CPU 执行时间的核心公式}{cpu-time-formula}
\textbf{CPU 执行时间}是衡量计算机性能的最终标尺。核心公式：
\[
\boxed{T_{\text{CPU}} = N \times \text{CPI} \times T_{\text{clk}} = \frac{N \times \text{CPI}}{f}}
\]
其中：
\begin{itemize}
    \item $N$：程序执行的总\textbf{指令条数}（Instruction Count）。
    \item $\text{CPI}$：平均每条指令的\textbf{时钟周期数}（Cycles Per Instruction）。
    \item $T_{\text{clk}}$：\textbf{时钟周期}（Clock Cycle Time），$T_{\text{clk}} = 1 / f$。
    \item $f$：主频（Clock Frequency）。
\end{itemize}

\textbf{三个维度的优化方向：}
\begin{itemize}
    \item 减少 $N$：编译器优化、更好的算法；
    \item 降低 CPI：微架构改进（流水线、超标量、乱序执行）；
    \item 减少 $T_{\text{clk}}$（提高主频）：改进工艺、更好的电路设计。
\end{itemize}
\end{formula}

\begin{conclusion}{MIPS 与 MFLOPS 计算}{mips-mflops}
\[
\boxed{\text{MIPS} = \frac{\text{指令条数}}{\text{执行时间} \times 10^{6}} = \frac{f}{\text{CPI} \times 10^{6}}}
\]
\[
\boxed{\text{MFLOPS} = \frac{\text{浮点运算次数}}{\text{执行时间} \times 10^{6}}}
\]

\textbf{注意事项：}
\begin{itemize}
    \item MIPS 是\textbf{与程序有关的指标}——同一 CPU 运行不同程序时 MIPS 不同，不能用 MIPS 直接比较不同指令系统的机器的性能。
    \item MIPS 忽略了指令功能的差异（复杂指令 vs 简单指令执行的工作量不同），因此有时会出现「MIPS 较高但实际性能却较差」的情况（所谓「MIPS 陷阱」）。
    \item MFLOPS 更适用于科学计算领域的性能比较。
\end{itemize}
\end{conclusion}

\begin{attention}
\textbf{408 高频计算考点——CPI 与 MIPS 的换算：}
\begin{enumerate}
    \item \textbf{已知主频和 MIPS，求 CPI：}
    \[
    \text{CPI} = \frac{f}{\text{MIPS} \times 10^{6}}
    \]
    \item \textbf{多类指令的加权 CPI：} 若程序包含 $k$ 类指令，每类指令条数为 $N_i$，CPI 为 $\text{CPI}_i$，则：
    \[
    \text{CPI}_{\text{avg}} = \frac{\sum_{i=1}^{k} N_i \times \text{CPI}_i}{\sum_{i=1}^{k} N_i} = \sum_{i=1}^{k} \left( \frac{N_i}{\sum N_i} \right) \times \text{CPI}_i
    \]
    \item \textbf{执行时间比较：} 比较两台机器的性能时，最终依据是 $T_{\text{CPU}}$（或与其成正比的 $N \times \text{CPI} / f$），不能只比较 MIPS 或主频。
\end{enumerate}
\end{attention}

\begin{conclusion}{性能指标速查表}{perf-metrics-table}
\begin{tablebox}{常见性能指标}
\centering
\begin{tabular}{|c|l|l|}
\hline
指标 & 定义 & 备注 \\
\hline
机器字长 & CPU 一次处理的二进制位数 & 32/64 位 \\
主频 $f$ & 时钟信号频率 & GHz 级 \\
$T_{\text{clk}}$ & $1/f$，时钟周期 & 例如 3 GHz $\to$ 0.33 ns \\
CPI & 每条指令的平均时钟数 & 越小越好 \\
IPC & $1/\text{CPI}$，每周期指令数 & 越大越好 \\
MIPS & 每秒百万条指令 & 与程序相关 \\
MFLOPS & 每秒百万次浮点运算 & 科学计算常用 \\
$T_{\text{CPU}}$ & $N \times \text{CPI} \times T_{\text{clk}}$ & 终极性能指标 \\
\hline
\end{tabular}
\end{tablebox}

其中 $\text{IPC} = 1/\text{CPI}$。现代超标量处理器的 IPC 可以大于 1（每个周期发射多条指令）。
\end{conclusion}

\begin{conclusion}{性能计算实例}{perf-example}
\boxed{\textbf{例}} 某计算机主频为 2.0 GHz，运行某程序共执行 $2\times 10^{9}$ 条指令，CPI = 1.5。求 CPU 执行时间和 MIPS。

\textbf{解：}
\begin{itemize}
    \item $T_{\text{clk}} = 1 / (2.0\times 10^{9}) = 0.5 \times 10^{-9} \text{ s} = 0.5 \text{ ns}$
    \item $T_{\text{CPU}} = N \times \text{CPI} \times T_{\text{clk}} = 2\times 10^{9} \times 1.5 \times 0.5 \times 10^{-9} = 1.5 \text{ s}$
    \item $\text{MIPS} = f / (\text{CPI} \times 10^{6}) = 2\times 10^{9} / (1.5 \times 10^{6}) \approx 1333.3$
\end{itemize}

\boxed{\textbf{例}} 某程序包含三类指令，A 类 10,000 条（CPI = 1）、B 类 20,000 条（CPI = 2）、C 类 10,000 条（CPI = 3）。求该程序的平均 CPI。

\textbf{解：}
\[
\text{CPI}_{\text{avg}} = \frac{10,000\times1 + 20,000\times2 + 10,000\times3}{10,000+20,000+10,000} = \frac{80,000}{40,000} = 2.0
\]
\end{conclusion}

\subsection{408 高频考点与答题策略}

\begin{conclusion}{本章 408 核心考点总结}{chap1-keypoints}
\begin{enumerate}
    \item \textbf{冯·诺依曼结构的核心特征}：存储程序、五大部件、以运算器为中心（注意与现代以存储器为中心的对比）。
    \item \textbf{五大部件的功能与辨析}：运算器（ALU + 寄存器）、控制器（PC/IR/ID）、两者合称 CPU；主机 = CPU + 主存。
    \item \textbf{软件分类}：系统软件（OS/编译器/汇编器/DBMS）vs 应用软件——这个知识点在选择题中常以「以下哪个属于系统软件」的形式出现。
    \item \textbf{三个语言级别}：机器语言（二进制）、汇编语言（助记符$\to$汇编器翻译）、高级语言（编译器翻译或解释执行）。注意区分编译与解释。
    \item \textbf{计算机层次结构}：$L_0$（微程序）到 $L_5$（应用层），理解每一层的作用及层次间的翻译/解释关系。常考点：哪几层是虚拟机？哪几层是物理机？
    \item \textbf{CPU 执行时间公式}：$T_{\text{CPU}} = N \times \text{CPI} \times T_{\text{clk}}$。这是串起所有性能指标的枢纽公式，必须熟练运用。
    \item \textbf{MIPS 计算}：$\text{MIPS} = f / (\text{CPI} \times 10^{6})$，注意单位的统一（主频用 Hz）。
    \item \textbf{多类指令的加权 CPI}：按频度加权求和，408 计算题高频出现。
\end{enumerate}
\end{conclusion}

\begin{attention}
\textbf{常见易错点：}
\begin{itemize}
    \item 主频单位换算：1 GHz = $10^{9}$ Hz；1 ms = $10^{-3}$ s；1 $\mu$s = $10^{-6}$ s；1 ns = $10^{-9}$ s。
    \item MIPS 公式中分母的 $10^6$ 不要遗漏。
    \item 冯·诺依曼计算机中「指令和数据放在同一存储器」——这里的「同一存储器」指的是逻辑上的统一编址主存，不是指 Cache 的物理分离。
    \item 「五大部分」中不包括「总线」——总线和电源属于支撑结构。
\end{itemize}
\end{attention}

\newpage
\section{习题}

\subsection{基础过关题}

\begin{enumerate}

\item \textbf{（单项选择题）} 冯·诺依曼计算机中，指令和数据在存储器中的存放方式是
\begin{center}
\begin{tabular}{ll}
(A) 分别存放在两个独立存储器中 & (B) 以同等地位存放在同一个存储器中 \\[4pt]
(C) 指令存放在 ROM，数据存放在 RAM & (D) 指令存放在 Cache，数据存放在主存
\end{tabular}
\end{center}

\item \textbf{（单项选择题）} 以下关于 CPU 组成的说法中，正确的是
\begin{center}
\begin{tabular}{ll}
(A) CPU 仅包含运算器 & (B) CPU 仅包含控制器 \\[4pt]
(C) CPU 包含运算器和控制器 & (D) CPU 包含运算器、控制器和主存储器
\end{tabular}
\end{center}

\item \textbf{（单项选择题）} 某计算机主频为 1.2 GHz，运行某程序的 CPI 为 2.0，则该计算机的 MIPS 为
\begin{center}
\begin{tabular}{ll}
(A) $600$ & (B) $1200$ \\[6pt]
(C) $2400$ & (D) $600\times 10^{6}$
\end{tabular}
\end{center}

\item \textbf{（单项选择题）} 以下不属于系统软件的是
\begin{center}
\begin{tabular}{ll}
(A) 操作系统 & (B) 编译器 \\[4pt]
(C) 数据库管理系统 & (D) 浏览器
\end{tabular}
\end{center}

\item \textbf{（填空题）} 在计算机系统的层次结构中，能直接被硬件识别和执行的层次是 \underline{\qquad\qquad} 级。汇编语言程序需要通过 \underline{\qquad\qquad} 翻译为机器语言。

\item \textbf{（简答题）} 简述冯·诺依曼结构的主要特点，并说明现代计算机结构与其最主要的区别是什么。

\item \textbf{（计算题）} 某计算机主频为 2.5 GHz，运行基准测试程序中共执行了 $5\times 10^{9}$ 条指令，CPI 为 1.6。求：（1）CPU 执行时间；（2）该计算机的 MIPS。

\end{enumerate}

\subsection{真题风格与综合训练}

\begin{enumerate}[resume]

\item \textbf{（真题风格，单项选择题）} 设某程序在计算机 A 上运行需 10 s，A 的主频为 2 GHz。设计算机 B 的目标是运行同一程序只需 6 s，且 B 的 CPI 为 A 的 1.2 倍，时钟周期为 A 的 0.9 倍，则 B 所需的主频至少为
\begin{center}
\begin{tabular}{ll}
(A) $2.56$ GHz & (B) $3.33$ GHz \\[6pt]
(C) $4.00$ GHz & (D) $4.44$ GHz
\end{tabular}
\end{center}

\item \textbf{（真题风格，单项选择题）} 某程序在某台计算机上运行，该程序包含 A、B、C 三类指令，其 CPI 分别为 1、2、3。已知程序中三类指令的条数之比为 $2:3:5$，则该程序的平均 CPI 为
\begin{center}
\begin{tabular}{ll}
(A) $2.0$ & (B) $2.1$ \\[6pt]
(C) $2.3$ & (D) $2.5$
\end{tabular}
\end{center}

\item \textbf{（真题风格，综合题）} 某计算机的主频为 3.0 GHz，运行某程序时测得 MIPS 为 1200。回答以下问题：
\begin{enumerate}
    \item[(1)] 求该程序运行时的平均 CPI；
    \item[(2)] 若该程序共执行 $3.6\times 10^{9}$ 条指令，求 CPU 执行时间；
    \item[(3)] 若通过编译器优化，将程序的指令条数减少 20\%，但 CPI 增加了 10\%，优化后的程序执行时间比优化前更快还是更慢？请给出计算过程。
\end{enumerate}

\item \textbf{（真题风格，综合题）} 在计算机系统的层次结构中：
\begin{enumerate}
    \item[(1)] 描述从 $L_4$（高级语言级）到 $L_1$（机器语言级）的转换过程，指出哪些步骤使用了翻译，哪些步骤使用了（或可能使用）解释；
    \item[(2)] 说明 Java 语言程序的执行过程如何体现「虚拟机」的概念；
    \item[(3)] 为什么说操作系统层（$L_2$）是虚拟机而非物理机？操作系统为上层提供了哪些抽象？
\end{enumerate}

\item \textbf{（真题风格，综合题）} 某处理器包含 4 类指令，其 CPI 和在某个典型程序中的使用频度如下表：
\begin{tablebox}{指令分布}
\centering
\begin{tabular}{|c|c|c|}
\hline
指令类型 & CPI & 频度 \\
\hline
整数 ALU & 1 & 45\% \\
数据传送 & 2 & 30\% \\
浮点运算 & 4 & 15\% \\
分支跳转 & 3 & 10\% \\
\hline
\end{tabular}
\end{tablebox}

\begin{enumerate}
    \item[(1)] 计算该程序的平均 CPI；
    \item[(2)] 若该处理器的主频为 2.0 GHz，求 MIPS；
    \item[(3)] 设计者计划将浮点运算的 CPI 从 4 降为 2，同时主频提升到 2.5 GHz，其他不变。改进后的整体性能提升为原来的多少倍？（用 $T_{\text{CPU}}$ 之比表示）
\end{enumerate}

\end{enumerate}

\vspace{8pt}
\noindent\textbf{拓展阅读：}
\begin{itemize}
    \item Patterson \& Hennessy《计算机组成与设计：硬件/软件接口》第 1 章「计算机概念与技术」——1.1--1.6 节全面覆盖本章内容，包括性能公式的导数和 Amdahl 定律。
    \item Bryant \& O'Hallaron《深入理解计算机系统》（CS:APP）第 1 章「计算机系统漫游」——1.1--1.9 节从程序员视角理解计算机系统的抽象层次，包含 Amdahl 定律（第 1.9 节）和并发/并行的初步讨论。
\end{itemize}
